%%%%%%%%%%%%%%%%%%%%%%%%%%%%%%%%%%%%%%%%%%%%%%%%%%%%%%%%%%%%%%%%%%%%%%
%%  NMH 论文 v3 中文对照检查版
%%
%%  本文件为 nmh-manuscript-v3.tex 的中文对照检查版（镜像翻译），
%%  生成日期：2026-09-27。
%%  用途：供作者审读修改中文内容；所有待填占位符以红色 \ph{...} 显示；
%%        本文件不用于投稿，投稿请使用英文版
%%        sn-article-template/nmh-manuscript-v3.tex。
%%  编译：latexmk -xelatex nmh-manuscript-v3-zh.tex
%%        （参考文献：\bibliographystyle{unsrt} + nmh-references-v3.bib）
%%
%%  占位符约定（与英文版一一对应）：
%%    [ModelName]、[Author N]、[N]、[X.X]、[X.XX]、[XX]、[4]、[8]、
%%    [Evidence needed: ...]、[Repository] 等 -> 均以红色 \ph{...} 呈现，
%%    其中纯数值占位符保持原样（如 [X.X] -> \ph{X.X}）。
%%
%%  v3 相对 v2 的修订要点（对应英文版文件头注释，供审读参照）：
%%    P0-1 结果部分五个结论式小节标题（其中依赖 T4 的持久性半句
%%        已降级为急性期措辞）
%%    P0-2 结果部分新增外部效应量基准对标段落
%%    P0-3 七组条件定义主表（G9 以表注给出）
%%    无标题的 6 段式引言，含明确的研究空缺句、与 Rollwage / Yosef
%%        的划界、以及对 Goldenberg 的回应
%%    任务式方法学小标题 + 独立"统计分析"节 + "报告摘要"；
%%        结论独立成节
%%    calibre 统一：12 次 CBT 会谈；评估时点 T0/S2..S12
%%        （长量表 K=10 于 T0/S12；中间时点用 5 条目短量表 K=5）；
%%        T4 为计划中的随访协议（不以现在时态报告数据）；
%%        严重度为按运行切换的配置维度（当前默认中度）；
%%        G7 阻断主要通用记忆写入
%%%%%%%%%%%%%%%%%%%%%%%%%%%%%%%%%%%%%%%%%%%%%%%%%%%%%%%%%%%%%%%%%%%%%%

\documentclass[11pt,a4paper]{article}

\usepackage[UTF8,heading=false,fontset=windows]{ctex}
\usepackage{geometry}
\geometry{margin=2.5cm}
\usepackage{amsmath,amssymb}
\usepackage{booktabs}
\usepackage{graphicx}
\usepackage{xcolor}
\usepackage{hyperref}
\hypersetup{colorlinks=true,linkcolor=blue!60!black,citecolor=blue!60!black,urlcolor=blue!60!black}

% 行距与段间距
\linespread{1.35}
\setlength{\parskip}{0.4em}
\setlength{\parindent}{2em}

% 自定义标题格式（蓝色章节标题）
\usepackage{titlesec}
\titleformat{\section}{\Large\bfseries\color{blue!40!black}}{\thesection}{1em}{}
\titleformat{\subsection}{\large\bfseries\color{blue!30!black}}{\thesubsection}{1em}{}

% 占位符高亮（红色），便于审读时定位
\newcommand{\ph}[1]{\textcolor{red}{[#1]}}

\title{\textbf{基于生成式智能体的抑郁症仿真治疗中\\认知重构与记忆巩固的 \emph{in silico} 分离研究}}
\author{NMH 稿件 v3 中文对照检查版 \\ \small 生成日期：2026-09-27}
\date{}

\begin{document}
\maketitle

\begin{center}
\rule{\textwidth}{0.5pt}\\[0.3em]
\small\textit{英文投稿标题：In silico separation of cognitive restructuring and memory consolidation\\
in simulated depression treatment with generative agents}\\[0.2em]
\small 短标题（running head）：Generative-agent separation of CBT mechanisms in silico\\[0.4em]
\small 作者：\ph{Author 1}~\ph{Surname}$^{1,*}$，\ph{Author 2}~\ph{Surname}$^{1}$，\ph{Author 3}~\ph{Surname}$^{1}$，\ph{Author 4}~\ph{Surname}$^{1}$\\
\small $^{1}$\ph{Department}，\ph{Institution}，\ph{City}，\ph{Country}\\
\small $^{*}$通讯作者：\ph{author1@example.org}；\ph{author2@example.org}；\ph{author3@example.org}；\ph{author4@example.org}\\[0.3em]
\rule{\textwidth}{0.5pt}
\end{center}

\vspace{0.5em}

%%====================================================================
\section*{摘要}
%%====================================================================

抑郁症是全球致残的主要原因之一，认知行为疗法（CBT）是一线治疗方案，然而 CBT 特定技术相对于非特异性治疗因素的贡献难以分离，且记忆巩固在疗效持久性中的作用无法在临床试验中被操纵。在此，我们提出 \ph{ModelName}——一个生成式智能体平台，通过阶段性主诉图状态机、逐轮情绪推断和四层提示词架构，将 Beck 认知模型操作化地嵌入虚拟社区。在一个带补充性积极接触条件的七组 \emph{in silico} 设计中，结构化 CBT 产生的 PHQ-9 降幅比频率匹配的支持性咨询多 \ph{X.X} 分（Cohen's $d = \ph{X.XX}$），由此分离出认知重构这一有效成分；记忆写入消融保留了急性期反应，持久性问题则保留给计划中的随访协议。仿真效应量已与已发表的荟萃分析范围进行基准对标 \ph{待补充证据：基准对照表}。该平台在 \emph{in silico} 中将特异性成分与非特异性成分分离开来；临床验证仍然必要。

\vspace{0.5em}
\noindent\textbf{关键词：}生成式智能体；抑郁症；认知行为疗法；记忆巩固；计算精神病学；大语言模型

%%====================================================================
%%  引言（无标题，6 段 Kambeitz 式漏斗结构：疾病负担 ->
%%  特异性/非特异性之争 -> 两个临床无法回答的问题 ->
%%  生成式智能体范式 + 效度立场 -> 研究空缺句
%%  （含与 Rollwage/Yosef 的划界）-> 本研究）
%%====================================================================

抑郁症是全球负担最重的精神障碍之一，影响约 2.8 亿人，预计到 2030 年将成为全球疾病负担的首要单一贡献者~\cite{who2022}。认知行为疗法（CBT）是国际临床指南普遍推荐的一线循证治疗，荟萃分析显示其相对于等待列表和常规治疗对照具有中到大效应量~\cite{cuijpers2013,hofmann2012}。然而，治疗结果存在异质性——约 40--60\% 的患者达到临床显著改善——且治疗结束后维持或侵蚀疗效收益的过程仍未被完全理解。

一个长期存在的争论是：CBT 的效应有多少可归因于其特定技术——认知重构、行为实验和结构化会谈推进——又有多少可归因于治疗联盟、定期专业接触和共情倾听等非特异性因素~\cite{lambert1992}。成分分析试验将完整方案与保留共同因子的拆解变体进行比较，提供了最直接的临床证据，并提示认知成分在情感支持之外贡献了额外获益~\cite{jacobson1996}。此类设计仍然罕见且受伦理约束：不能为研究目的而拒绝向真实患者提供已确立治疗中的任何活性成分，因此特异性与非特异性的问题从未在析因设计所能提供的受控分离条件下得到回答。

两个机制层面的问题在临床上尤其难以回答。第一，分离 CBT 技术的特异性效应需要一个在不使用这些技术的前提下提供相同非特异性因素的对照条件——这一设计在伦理和实践中都难以在真实患者中实施，因为不能拒绝向患者提供循证的活性治疗。第二，检验记忆巩固——即治疗性经验的内部编码与复述——是否因果性地维持治疗效果的持久性，需要操纵记忆形成过程，而除药物或基于睡眠的范式（其自身带有混淆因素）之外，这在人类身上无法实验性地实现~\cite{rasch2013}。尽管理论论述将持久改变与认知--情感表征的反复激活和再巩固联系起来~\cite{teasdale1988,spens2024}，但在患者身上对技术、记忆和症状动态的过程层面观察，仍局限于回顾性自我报告和间歇性评估。

生成式智能体——由大语言模型（LLM）驱动、在虚拟社区中维护记忆流、进行规划、反思和对话的拟真模拟体~\cite{park2023}——提供了一种互补范式，且基于智能体的 LLM 仿真正在社会科学中迅速扩散~\cite{gao2024}。Park 等人引入了一个能大规模产生可信社会行为的框架~\cite{park2023}；Kambeitz 和 Meyer-Lindenberg 提出此类智能体可用于建模环境和社会决定因素对心理健康的影响~\cite{kambeitz2025npj}。LLM 角色扮演可作系统性分析~\cite{shanahan2023}；人格特质可通过心理测量工具在 LLM 中被测量和塑造~\cite{serapio2025,wright2026}；先复现已知规律、再运行反事实情境的仿真研究为保真性推理提供了模板~\cite{qu2024,ji2026}。由于大语言模型并不具备情绪~\cite{goldenberg2026}，此类平台的效度主张不能建立在现象体验之上；可辩护的目标是行为保真性——即模拟患者能否再现抑郁症的言语、活动和量表反应规律，这是从计算社会科学改编而来的算法保真性准则~\cite{qu2024}。

心理健康研究已开始占据相邻领域：数字认知孪生~\cite{doraiswamy2025}、聊天机器人与用户信念之间的反馈回路~\cite{dohany2026}以及计算精神病学的社会智能体扩展~\cite{rhoads2024}都是活跃话题；精神病理学的实证结构已在 LLM 内部参照真实临床数据得到探测~\cite{kambeitz2025nmh}；LLM 架构已作为向真实被试传递认知行为技术的治疗师得到评估~\cite{rollwage2026}；微调的治疗师模型已由数字患者在单次会谈交互中加以评估~\cite{yosef2025}。然而，迄今为止尚无研究将临床级别的抑郁症精神病理学嵌入生成式智能体社区，并使结构化心理治疗协议在纵向进程中接受受控的多臂 \emph{in silico} 评估：现有的抑郁症模拟器依赖在脱离任何社会情境的情况下对孤立患者模型进行专家标注的微调~\cite{liu2025}，而 \emph{in silico} 治疗比较迄今使用的是生物物理简化模型，而非互动的患者与治疗师~\cite{montague2012,hauser2022,deco2024}。

在此，我们提出 \ph{ModelName}——一个生成式智能体平台，通过阶段性主诉图状态机、带波动约束的逐轮情绪推断和四层提示词架构，将 Beck 抑郁认知模型~\cite{beck1967}操作化地实现，并整合 12 次会谈的结构化 CBT 协议和阶段性心理测量评估组合。我们报告一项七组对照实验——包括频率匹配的支持性咨询（G6）和保留完整 CBT 的记忆写入消融（G7），以及补充性积极接触条件（G9）——旨在分离两个长期被混淆的治疗成分：认知重构与非特异性治疗因素，以及记忆写入在持久性中的仿真内作用，其中后半部分保留给计划中的随访协议。我们既将该平台呈现为一个对照已知规律检验真实性的抑郁状态模型，也将其呈现为一个双机制分离，同时承认 \emph{in silico} 因果主张仅指仿真的记忆写入机制，需要临床转化。

%%====================================================================
%%  结果（五个结论式小节标题；第四个标题中依赖 T4 的持久性半句
%%  已降级为急性期措辞；第五个标题采用架构论证版本）
%%====================================================================
\section{结果}\label{sec:results}

\subsection*{生成式智能体再现了与严重度一致且内部连贯的抑郁状态}

我们从五个维度评估抑郁状态的真实性：状态连续性（主诉阶段之间无不合理跳变）、情境一致性（同一主诉在工作、家庭和治疗情境中以不同方式表达）、关系差异性（在医生、朋友和家人面前调整信任度、防御性和自我袒露）、主诉可解释性（核心信念和叙事焦点可在话语中追溯）以及言语风格稳定性（语速、袒露水平、语气）。在阶段性评估协议下，\ph{XX}\% 的轨迹片段满足这些描述性真实性标准，其中情境一致性通过率为 \ph{XX}\%，关系差异性通过率为 \ph{XX}\% \ph{待补充证据：描述性真实性编码；专家盲评协议尚未就绪}。作为基于量表的检验，在 T0 施测的基线 PHQ-9~\cite{kroenke2001}（每个冻结快照 $K=10$ 次重复生成）与所分配的严重度档位一致——当前中度默认批次中为中度 \ph{X.X$\pm$X.X}，早期重度批次中为重度 \ph{X.X$\pm$X.X}（逐批次标注）——且配置了较高神经质的人设产生了更高的基线 PHQ-9（均值差异 \ph{X.X}，Cohen's $d = \ph{X.XX}$），在方向上复制了神经质--抑郁关联，构成第二种形式的外部锚定（探索性；人设覆盖有限）。冻结快照内的重复生成信度为 $\mathrm{ICC}(1,K) = \ph{X.XX}$（扩展数据图~2）\ph{待补充证据：重复汇总计算}，且 PHQ-9 与 BDI-II~\cite{beck1996} 总分在同期施测长量表的评估节点上相互趋同（$r = \ph{X.XX}$）。一套拟人化验证组合——会谈层面的话语状态与情绪标注、锚定个案参数的静态画像恢复，以及对主诉图变化证据的独立盲审——已在方法部分作出规定，尚待目标批次计算 \ph{待补充证据：PSI-Bench 式指标；人类参照数据尚未获得}。系统架构总结于图~\ref{fig:architecture}，人设特征见扩展数据表~1。

\subsection*{\emph{in silico} 治疗反应按干预类型与社会接触效价呈梯度分布}

在每条件 \ph{N} 次独立运行中，结构化 CBT（G1）产生了从 T0 到 S12 平均 \ph{X.X} 分的 PHQ-9 降幅（较基线 \ph{XX}\%，Cohen's $d = \ph{X.XX}$，95\% CI \ph{X.X, X.X}），且 BDI-II 在 S12 与 PHQ-9 趋同（$r = \ph{X.XX}$）。在当前 72 步协议下，PHQ-9 仅在 T0 和 S12 施测，中间时点（S2--S10）施测五条目抑郁短量表；在该短量表负荷上，改善最陡峭的区段落在第 \ph{4} 至第 \ph{8} 次会谈之间 \ph{待补充证据：当前协议下短量表轨迹的分段与阶段归因；归因将基于固定的 12 次会谈大纲重新推导}。无干预组（G2）显示 \ph{X.X} 分的变化（\ph{XX}\%，与自然病程一致）。社会接触条件形成了梯度模式：积极社会暴露（G9）$>$ 中性社会接触（G3）$>$ 无干预（G2）$>$ 负面社会暴露（G5，使症状恶化 \ph{X.X} 分），支持了不良社会环境主动加剧而非仅仅未能改善症状的假设；G3 这一环节有待同版本重跑 \ph{待补充证据：同版本 G3 运行}。咨询室条件（G4）与 G1 相差 \ph{X.X} 分（$d = \ph{X.XX}$，$p = \ph{X.X}$）；由于 G4 在继承完全相同的医患干预的同时，把完整小镇替换为双角色咨询室场景，这一差异标定的是场景设置的综合效应，而非单一的环境受限效应。为给这些量级提供外部锚定，我们将仿真的 G1 效应量逐点对标已发表的心理治疗荟萃分析范围~\cite{cuijpers2013,hofmann2012,cuijpers2020}，以及对话式智能体递送 CBT 的随机对照与荟萃分析证据~\cite{heinz2025,hang2026}：仿真效应（扩展数据表~3）相对这些范围落于 \ph{待补充证据：效应量基准对照表；需单一锁定批次的校准}。如有偏离，将连同区间与方向透明报告，不做任何事后协议切换；若偏离较大，该基准对标将降格为校准性讨论，而非等效性主张。纵向轨迹见图~\ref{fig:trajectories}；条件定义见表~\ref{tab:conditions}；基线与终点统计见扩展数据表~2。

\subsection*{认知重构在非特异性治疗因素之外贡献了额外获益}

关键的 G1 对 G6 比较分离了 CBT 技术的特异性贡献，把成分分析传统~\cite{jacobson1996}移植到一个完全受控的 \emph{in silico} 环境中：两组接受相同频率的医患接触和相同的咨询历史支持，但 G6 使用非指导性、以确认为中心的沟通，不含认知重构、行为实验或结构化会谈推进。结构化 CBT 比单纯支持性咨询多产生了 \ph{X.X} 分的 PHQ-9 降幅（Cohen's $d = \ph{X.XX}$，95\% CI \ph{X.X, X.X}，bootstrap），这一优势在现有人设集上方向一致。鉴于现有批次中每条件独立运行数量较少，我们强调带置信区间的描述性效应量而非推断性检验 \ph{待补充证据：确认性检验之前各批次的外层运行覆盖}。这表明认知重构——歪曲识别、证据检验和行为实验——在联盟、定期接触和共情倾听之上贡献了治疗获益。终点分布和两两效应量见图~\ref{fig:endpoints}。

\subsection*{阻断记忆写入保留了急性期治疗反应}

记忆写入消融（G7）接受了与 G1 完全相同的结构化 CBT，但阻断了患者智能体的主要通用记忆写入——事件、思想和聊天节点——同时动态抑郁状态机继续运行。在急性终点（S12），G7 的 PHQ-9 与 G1 无差异（$d = \ph{X.XX}$，$p = \ph{X.X}$），表明在没有新记忆形成的情况下急性期治疗反应得以保留（图~\ref{fig:cmech}）；现有 G7 存档来自早期方法版本，因此该比较带有明确的版本边界，需待同版本重跑后方可作推断使用 \ph{待补充证据：同版本 G7 运行}。阻断记忆写入是否损害持久性，保留给计划中的随访协议（T4）：第 12 次会谈之后是一个无医患接触的独立期间，期间不施测任何量表以避免扰动自然轨迹，期末从冻结检查点恢复终点分数。随访节点仅在早期 120 步批次中演练过，不属于当前 72 步配置的一部分；因此 T4 仅作为计划中、待执行的内容报告，而非已获得的结果。如果在 T4 观察到"急性保留而持久性受损"的分离，则与记忆写入在维持治疗收益中的仿真内因果作用相一致~\cite{teasdale1988,spens2024}；这是否推广到人类记忆巩固，是一个需要临床转化的推论。在存档版本中，治疗前后的主诉图阶段性转移在 G1 与 G7 之间存在差异，G1 向残余阶段推进的幅度更大（扩展数据图~1）\ph{待补充证据：同版本主诉图证据，按控制器分层}。

\subsection*{整合式认知架构是真实状态动态所必需的}

上文报告的动态并非来自单一的端到端生成器，而是来自架构组件之间的交互；在此我们将架构层面的论证与结局层面的论证并举。阶段性主诉图为重构提供其运作的基底：治疗对话以阶段锚定的歪曲认知为目标，保真性门控的阶段更新协议将长期轨迹约束在临床可解释的区域内，这正是上文状态连续性与主诉可解释性得以通过、而非出现自由漂移的原因。逐轮情绪层通过波动约束以及明确禁止把瞬态回读为持续改善，防止不合理的症状跳变并支持轨迹连贯性。四层提示词栈使同一底层状态在治疗、工作和家庭情境中呈现为不同形态，这是情境一致性与关系差异性背后的机制。这些论证属于设计层面且跨条件相关；决定性检验是模块级消融——NoGraph、NoEmotion 和 NoRelation 臂——尚未执行，相应运行已列入计划 \ph{待补充证据：模块级消融实验，约需额外 9 次运行；各臂定义超出当前方法学规范，属新增内容，待方法学侧批准}。在这些运行完成之前，本小节的主张仍是架构层面的论证，最终将由实验结果取代。

%%====================================================================
%%  讨论（按 NMH 政策不设小节）
%%====================================================================
\section{讨论}\label{sec:discussion}

\ph{ModelName} 表明，一个嵌入 Beck 认知模型的生成式智能体平台能够产生内部连贯、与严重度一致的抑郁状态动态，并通过受控的 \emph{in silico} 设计开始分离 CBT 中两个长期被混淆的成分：认知重构（通过 G1 对 G6 分离）和记忆写入在持久性中的仿真内作用（通过 G7 探究，持久性部分保留给计划中的随访协议）。观察到的效应量级与结果部分报告的外部基准对标相呼应；鉴于仿真的自我报告评估，这是量级参照而非等效。这一双重框架把两个问题升级为同一个框架：治疗反应具有振幅成分和持久性成分，由不同机制决定，二者均可在 \emph{in silico} 中被分离。

这项工作把生成式智能体范式从一般社会模拟~\cite{park2023,gao2024}扩展到以临床为根基的心理健康研究，从概念性提案~\cite{kambeitz2025npj,doraiswamy2025,rhoads2024}走向受控实验，并与生物物理基础模型中的 \emph{in silico} 治疗比较~\cite{deco2024}以及关于社会推断和情绪的计算精神病学论述~\cite{barnby2024,trier2025,zavlis2025}形成互补。与参数化的基于智能体模型~\cite{bonabeau2002}相比，\ph{ModelName} 以自然语言——即真实心理治疗运作的媒介——呈现主诉图锚定的认知、逐轮情绪和治疗对话。其评估哲学遵循基于 LLM 的社会仿真"先复现、后反事实"的逻辑~\cite{qu2024,ji2026}。与此同时，对话式临床 AI 已在盲评协议下完成大规模评估~\cite{tu2025,saab2026,johri2025}，包括作为心理治疗智能体~\cite{rollwage2026}以及作为人类同伴帮助者的支持~\cite{sharma2023}，关于 AI 递送治疗的证据正在积累~\cite{heinz2025,hang2026,hua2025,stade2024}。我们的贡献是互补性的：智能体充当实验被试而非临床行为者；角色扮演是工具而非主张~\cite{shanahan2023}；人设以心理测量学方式处理~\cite{serapio2025,wright2026}；且优先采用受控多臂设计而非自由形式的多智能体协作~\cite{kim2026}。已部署聊天机器人的用户经验记录与信念动态风险~\cite{maples2024,siddals2024,dohany2026}进一步支持把此类系统保持为研究工具。

\emph{in silico} 优先的策略有伦理上的理由，而不仅仅是便利。在患者身上操纵记忆巩固是不可允许的，而为了量化非特异性因素而在七个臂中拒绝提供已确立的活性治疗，在临床上无法被证成；模拟被试不受此类约束，且合成数据治疗评估已有近期先例~\cite{yosef2025}。由于模拟对话中可能出现危机或自伤内容，独立的安全机制治理与风险相关的输出，且平台不进行诊断、风险评估或治疗决策。

本研究存在若干局限。第一，平台仅为研究目的模拟抑郁相关的表现和求助行为。第二，这里的 PHQ-9 和 BDI-II 是 LLM 模拟的自我报告，采用逐项无状态协议、混合规则评分和一致性校验，并非临床施测工具；因此与真实荟萃分析的效应量比较是量级参照而非等效。第三，G7 在动态抑郁状态机继续运行的同时阻断主要通用记忆写入——这是一种没有直接临床对应物的极端二值操纵，因为真实世界的记忆衰减是渐进的——且阻断的拦截覆盖率有待日志审计 \ph{待补充证据：基于存档运行日志的拦截覆盖率审计}。第四，人设和严重度覆盖相对于真实临床人群仍然有限；在更多档位重跑之前，当前的中度默认设置限制了按严重度分层的主张。第五，\emph{in silico} 因果主张指仿真的记忆写入机制——即仿真内因果——而非人类记忆巩固，后者仍是一个有待临床转化的假设。

未来工作应扩展人设在年龄、性别和文化上的多样性；纳入药物治疗建模；将仿真参数与临床试验数据集进行校准；完成计划中的模块级消融和随访协议；并检验 \emph{in silico} 中观察到的记忆巩固分离能否预测临床队列中的复发模式。模块化架构支持患者--治疗匹配的组合探索，迈向 \emph{in silico} 精准精神病学。

%%====================================================================
%%  结论（独立成节，不超过 80 词）
%%====================================================================
\section{结论}\label{sec:conclusions}

\ph{ModelName} 提供了一个生成式智能体平台，其中结构化 CBT、频率匹配的支持性咨询和记忆形成可以被分别操纵，从而分离出认知重构这一有效成分，并探究记忆写入在治疗持久性中的仿真内作用。因果主张以仿真为边界；临床转化需要对照真实队列进行验证。

%%====================================================================
%%  方法（任务式小标题；统计单独成节；报告摘要置于最后）
%%====================================================================
\section{方法}\label{sec:methods}

\subsection*{生成式智能体仿真平台}

\ph{ModelName} 扩展了 Stanford Town 生成式智能体框架~\cite{park2023}。村庄模式实例化一个完整的虚拟社区；咨询室模式将场景缩减为两个角色（即 G4 条件）。干预条件是深度合并到基础配置之上的配置叠加层；批次管线分两层组织：先批量仿真，随后冻结快照并进行独立的重复复评。规则负责时间步进、地图寻路、会面调度、条件分配、会谈推进和快照触发；LLM 负责日程规划、自然语言交互、反思、情绪推断和部分会谈判断。当前协议推进 72 个仿真步（每步 720 分钟，共 36 个仿真日），每 6 步安排 1 次受控会谈，共 12 次；一次医生会谈持续 2,881 个仿真分钟，一次居民强制聊天持续 1,440 分钟。每个智能体维护由事件、聊天和思想节点构成的记忆流，按新近性、重要性和相关性以 BGE-M3 嵌入进行检索；每个人设初始化时带有四条压力记忆（一个压力事件、一个羞耻触发、一个消极自我念头和一个情境性紧张），咨询历史检索则汇总先前会谈以保持连续性。在 G1 链条下，每次会谈最多提取三个家庭作业，并交由一个环境模型处理，其结果作为记忆注回。一个受 PatientAct 式患者建模启发的动态抑郁状态机在所有组中独立于主诉图运行。模型分配（仅公开可写配置；不报告任何端点或密钥）：DeepSeek 服务于医生、患者和评估链条；本地 Qwen3-8B 服务于离线审查和本地思考 LLM 调用（会谈任务规划、候选选择、摘要压缩）；检索嵌入运行在独立部署上。仅医患对话启用思考模式。模型身份与推理选项属于运行配置，按批次清单报告，而非固定的方法学常数。

\subsection*{抑郁认知架构}

患者认知通过四层提示词栈呈现：(i) 基础人格层（姓名、年龄、特质、习惯、日程安排及当日偶发事件）；(ii) 当前主诉阶段层（阶段标签与摘要、近期相关经历，以及由根主诉锚、固定核心信念与表达风格指导构成的稳定背景）；(iii) 会话上下文层（地点、时间、交互对象及其关系、互动类型与场景级判断）；(iv) 逐轮情绪与表达指导层（情绪标签、风格、强度、0--10 量表上的袒露度与防御性、相对上一轮的波动幅度，以及自我纠正模式）——并受一项明确约束：瞬态状态不得被回读为持续的症状改善。渲染过程还并入咨询历史记忆、智能体的记忆列表与一个滚动对话窗口。第五个"认知偏差"层仅作为未启用的设计草稿存在，从不被渲染。主诉图更新遵循一个阶段化的三步协议：三次相互独立、带代码门控的 LLM 调用。变化检测器仅依据经授权的患者文本，判断相对当前阶段是否出现了足以改变后续表达或行为的最小变化；判断处于密集的反例约束之下（计划不等于执行；瞬态症状不等于持续的能力丧失；带限定语的陈述不得被截断为前半句）。阶段规划器将下一阶段提议为"当前阶段 + 经验证的变化"，标签不超过 100 字符，摘要不超过两句（一句写持续存在的主诉，一句写最小变化及其限度），父主诉轴保持不变，且不扩写疗效、机制或未来前景。独立的阶段验证器出具变化保真性与阶段保真性两项裁定，二者全部通过方可提交；验证器逐一核查每条新断言，防止创建同义节点，防止夸大改善或恶化。三次调用及其输入与结果全部留有追踪。边界由代码强制执行：个案背景、根主诉锚与核心信念在初始化时即冻结；风格、情绪基线与关系参数为继承所得，不能被 LLM 提议改写；新节点标识符与边仅在成功提交时分配；终末阶段不再接受后继节点；患者原话与审计记录不可篡改；医生话语、反思或 CBT 进展绝不充当患者事实的证据——只有正式的患者表达才计为新证据，证据不足时主诉图绝不推进。

\subsection*{干预与对照条件}

医生智能体递送一个具有固定大纲顺序（会谈 1，随后单元 2.1--2.3、3.1--3.3 和 4.1--4.4）的 12 次会谈结构化 CBT 协议，从签到开始，经由自动化思维识别、认知歪曲标记、苏格拉底式重构和行为激活，直到复发预防；共设计了六套治疗套件（CBT、IPT、BA、DBT、PST 与 SPT），CBT 为主链实现。会谈内部控制委托给本地组件：会谈任务规划器每轮单调更新微目标（从无到部分到完成）；选择器从 Router 控制的候选集中挑选回复策略与微技能，无需额外的大模型调用；轻量级对话判断器仅见长期个案背景与可观察的结构化状态（绝不接触患者内部变量），处理有限数量的判断与终止，终止仅在微目标进度阈值之后、或估计会谈时长约 40 分钟之后才被开放——更早的终止输出一律强制为 false。会话上下文是先前会谈摘要加最近五轮医患交流的滚动窗口，会谈时长按每分钟 180 字符估算。完成门控采用 60\% 进度阈值与 60 轮硬上限：自第 50 轮起转向以巩固为导向的收尾，自第 55 轮起给出有原则的收尾建议；三级风险 Router（无、可能、高）始终生效。居民干预由组配置中硬编码的三个固定居民角色递送，不受患者人设影响；受控调度条件每 6 步注入一次脚本前缀（第 2 阶段），含中性、负面与积极三种效价版本——中性对话限于日常事务；负面脚本结合低回应性、情感错位与挫折，同时排除辱骂、威胁或鼓励自伤；积极脚本提供简短、具体、可信的支持，不含泛泛的打气或建议。自然发生变体——仅在居民于未被提示的对话中主动与患者交谈时注入前缀——存在于更广的设计空间中，此处不予分析。表~\ref{tab:conditions} 定义了 G1--G7 条件，补充性积极接触条件 G9 在表注中定义。控制器身份或提示词摘要不同的批次分开分析，绝不合并。

\subsection*{阶段性心理测量评估}

评估时点为 T0、S2、S4、S6、S8、S10 和 S12，每两次会谈一个中间时点。T0 和 S12 施测 PHQ-9（九条目）~\cite{kroenke2001} 和 BDI-II（21 条目）~\cite{beck1996}，每个冻结快照各 $K=10$ 次独立重复生成；中间时点施测两个五条目短量表——一个总体抑郁程度与影响量表和一个总体焦虑程度与影响量表（每条目 0--4，总分 0--20）——各 $K=5$。没有安排会谈的组每 12 步获得虚拟对齐标签。仿真期间平台只冻结状态（仅捕获）；一个独立的存档脚本对快照进行复评，选取每次运行中实际实现的阶段性节点。T0 在初始压力记忆注入之后、患者第一个自主日常决策之前捕获。快照是可恢复的运行状态，携带运行配置、触发元数据、带本地存储副本的智能体检查点、对话与动态抑郁状态（含个案级核心信念）、内容哈希以及组/人设/严重度/控制器溯源。每个条目的作答不带记忆：每个条目新建一个单条目临时对话，不向聊天记忆追加内容，不做动态提交，条目之间互不继承——这一无状态协议不同于连续的临床访谈，并照实报告。长量表由一个专家 LLM 从自然语言回答评分（每条目 0--3，总分、严重度和安全字段）；短量表采用混合评分：规则优先提取无歧义的 0--4 条目，完全可规则评分的答卷完全跳过 LLM，无法解析的条目回退到 LLM 且规则评分覆盖 LLM 评分，全程处于评分来源审计之下。基于规则的校验器复查条目数（9/21/5）、条目取值范围、加和一致性、严重度分级以及 PHQ-9 第 9 条与安全字段的一致性。治疗后随访协议（T4）已列入计划：第 12 次会谈之后是一个无医患接触、不施测量表的独立期间，期末从冻结检查点恢复分数；随访节点仅在早期 120 步批次（FU30--FU120）中演练过，当前 72 步配置中没有随访节点。拟人化验证组合规定如下，计算待完成：对每次会谈中患者的前八条回复做单选话语状态标注（困扰、探索/反思、领悟/接纳、填充语）与九种基本情绪标注，仅使用其前文语境（话语状态用六条消息，情绪用四条）；针对三个锚定个案参数（根主诉锚、核心信念、长期背景），在三个干扰项中做静态画像恢复；以及一项离线的主诉图变化证据审计，对图更新、反思、量表和分组保持盲视，仅见患者话语与活动档案，引用可溯源 \ph{待补充证据：PSI-Bench 式指标；人类参照数据未获得}。

\subsection*{实验设计}

实验为七组对照比较（G1--G7）加一个补充性积极接触条件（G9），依照仿真研究惯例~\cite{crosland2025}组织为一张场景定义表。比较单位是独立仿真运行；冻结快照内的重复生成刻画的是测量不确定性，不计入独立 $N$。两套人设系统并存：其一是包含九个人设的临床情境案例库，每个人设配有三个运行时可切换的严重度档位（轻度、中度、重度）以及额外的未分档默认配置；其二是用于小镇实验的八人设选择器集，年龄覆盖 18--66 岁，具有中度核心信念与可恢复的活动（仅中度档）。严重度是一个按运行切换的配置维度：当前默认为中度（自 0908 批次起）；更早时期使用中/重度混合或统一重度，凡引用早期批次的结果均标注该批次所属档位。
%% 严重度档位备注（分歧已标记，待人工裁定）：方法学摘要（截至 0923 日志）把小插图库
%% 记录为 9 个人设 x 4 个档位（含一个未分档的"正常"档）= 36 种配置；当前检出的仓库
%% 仅显示三个运行时可切换档位（runshells/kabuda_variant_runtime.py:106-110
%% SEVERITY_CONFIG_NAMES、:321 --severity 选项；runshells/run_batch_experiment.py:198
%% SEVERITIES），且 docs/town_method/01_experimental_conditions.md:142 统计出 9x3=27 个
%% 选择器可覆盖的 KBD 条件，未分档的 depression_config.json 默认值位于 assets 中。
%% 此处措辞遵从仓库实测，待裁定。外层重复没有全局统一取值（视批次而定，每条件 1--5 次），
%% 按批次清单报告，连同控制器身份（legacy、minimal 或 progressive-D）一起，
%% 后者作为分层变量。
%% 译注：英文原稿中"Outer repeats have no global value ..."一句位于 %% 注释块内，
%% 未在英文编译稿中显示，故此处同样保持为注释。

\subsection*{统计分析}

主要终点是 T0 到 S12 的 PHQ-9 变化；纵向描述使用中间短量表序列（S2--S10）。计划中的主分析框架是线性混合效应模型（变化 $\sim$ 组别 $+$ (1$|$persona) $+$ (1$|$run)）；鉴于现有批次中每条件独立运行数量较少（每条件 \ph{N} 次），分析以描述性效应量为主——Cohen's $d$~\cite{cohen1988} 附 bootstrap 95\% 置信区间——辅以 Wilcoxon 秩和敏感性检验，并在 Bonferroni 校正之外进行假发现率控制，遵循对话式 AI 的评估惯例~\cite{tu2025}。测量信度采用快照内样本标准差、95\% $t$ 区间、重复生成的组内相关系数 ICC(1,1) 与 ICC(1,$K$)~\cite{shrout1979}，以及条目级二次加权 kappa；PHQ-9/BDI-II 的趋同性在同期施测长量表的节点上评估。分析脚本和统计输出在代码仓库中可用。

\subsection*{报告摘要}

有关研究设计的更多信息，见与本文关联的 Nature Portfolio 报告摘要。报告遵循行为科学中大语言模型的使用清单~\cite{feuerriegel2026}：模型身份、推理选项和控制器身份按批次清单披露；本研究为纯仿真，不涉及人类被试和个人数据；大语言模型（DeepSeek、Qwen3-8B）在负责任发展的考量下作为研究基础设施使用~\cite{stade2024}；未经作者审阅，任何 LLM 输出均未纳入稿件正文。

%%====================================================================
%%  图表（正文：4 图 + 1 张条件定义表；扩展数据：2 图 + 3 表）
%%====================================================================

\begin{figure}[h]
\centering
%% \includegraphics[width=0.9\textwidth]{figure1_architecture.pdf}
\fbox{\parbox[c][6cm][c]{0.9\textwidth}{\centering \emph{\ph{图 1 占位}}\\架构与流程概览：虚拟社区（村庄 / 咨询室）$\rightarrow$ 抑郁认知架构 $\rightarrow$ 带本地会谈控制的 12 次会谈 CBT 管线 $\rightarrow$ 阶段性评估与七组实验设计。}}
\caption{\ph{ModelName} 的系统架构与实验流程。生成式智能体虚拟社区（村庄或咨询室模式）与抑郁认知架构（四层提示词栈、阶段性主诉图、逐轮情绪推断、动态抑郁状态机）、带本地会谈控制的 12 次会谈结构化 CBT 管线（任务规划器、选择器、带延迟终止的轻量级对话判断器），以及阶段性评估框架（仅捕获的冻结快照加独立重复复评）相耦合，处于七组实验设计之下。}
\label{fig:architecture}
\end{figure}

\begin{figure}[h]
\centering
%% \includegraphics[width=0.9\textwidth]{figure2_trajectories.pdf}
\fbox{\parbox[c][6cm][c]{0.9\textwidth}{\centering \emph{\ph{图 2 占位}}\\各组纵向评估轨迹：S2--S10 的五条目抑郁短量表，T0 和 S12 的 PHQ-9 锚点，95\% CI 带。}}
\caption{各组纵向评估轨迹。中间时点 S2--S10 施测五条目抑郁短量表；PHQ-9 锚定 T0 和 S12（每个冻结快照各 $K=10$ 次重复生成）。G2 组标签为虚拟对齐标签。改善最陡峭的区段以短量表负荷表达 \ph{待补充证据：当前协议下的分段与阶段归因}。阴影带表示跨独立运行的 95\% 置信区间。}
\label{fig:trajectories}
\end{figure}

\begin{figure}[h]
\centering
%% \includegraphics[width=0.9\textwidth]{figure3_endpoints.pdf}
\fbox{\parbox[c][6cm][c]{0.9\textwidth}{\centering \emph{\ph{图 3 占位}}\\终点 PHQ-9 分布，附两两 Cohen's $d$（高亮 G1 vs G6）。}}
\caption{各组终点 PHQ-9 分布，附两两效应量。G1 对 G6 的对比把认知重构从非特异性治疗因素（频率匹配的接触、无技术成分）中分离出来。数值待目标批次核验。}
\label{fig:endpoints}
\end{figure}

\begin{figure}[h]
\centering
%% \includegraphics[width=0.9\textwidth]{figure4_cmech.pdf}
\fbox{\parbox[c][6cm][c]{0.9\textwidth}{\centering \emph{\ph{图 4 占位}}\\记忆写入消融：急性期 S12 终点（G7 vs G1）与计划中的 T4 随访设计（独立无接触期间、冻结检查点恢复）。}}
\caption{记忆写入消融把急性期反应与持久性分离开来。G7 接受与 G1 相同的 CBT，但阻断主要通用记忆写入（事件/思想/聊天；动态抑郁状态机继续运行）。急性期 S12 终点在明确的版本边界下与 G1 比较 \ph{待补充证据：同版本 G7 运行}；T4 随访协议——一个无中期量表、期末从冻结检查点恢复分数的独立无接触期间——已列入计划、待执行。}
\label{fig:cmech}
\end{figure}

\begin{table}[h]
\caption{实验条件定义（G1--G7；当前 72 步协议：一个仿真步 = 720 分钟，每 6 步安排一次会谈）。G9 是七组框架之外的补充条件——受控的积极居民聊天，按 G3 的方式调度、使用积极效价脚本——与 G2/G3/G5 一起构成社会接触效价梯度。G3 和 G7 行有待同版本重跑后作定量对比 \ph{待补充证据}；G4 行标定的是场景的综合差异，而非单一干预效应。}\label{tab:conditions}
\footnotesize
\setlength{\tabcolsep}{3pt}
\begin{tabular}{@{}p{0.9cm}p{1.8cm}p{2.0cm}p{2.8cm}p{2.3cm}p{2.0cm}@{}}
\toprule
组别 & 条件 & 接触与调度 & 技术内容 & 记忆状态 & 临床问题 \\
\midrule
G1 & 结构化 CBT（基准） & 每 6 步一次医生会谈；共 12 次 & 完整 12 次会谈 CBT 大纲，含会谈提示词、对话判断器、会谈评估、咨询历史及作业--环境回路 & 完整记忆写入；强制聊天记忆策略 & 完整 CBT 管线的疗效 \\
G2 & 无干预 & 无；量表按步触发并附虚拟会谈标签 & 无 & 完整记忆写入；动态抑郁激活 & 自然病程 \\
G3 & 受控中性居民聊天 & 每 6 步一次居民聊天（1,440 仿真分钟） & 中性日常事务脚本 & 完整记忆写入 & 有陪伴交谈与无陪伴之对比 \\
G4 & 咨询室场景设置 & 每 6 步一次医生会谈，同 G1 & 与 G1 相同的基础 CBT 医患干预，置于双角色场景 & 完整记忆写入 & 完整小镇与双角色咨询室的综合差异（场景消融） \\
G5 & 受控负面居民聊天 & 同 G3 & 负效价脚本（低回应性、情感错位、挫折） & 同 G3 & 不良社会暴露 \\
G6 & 支持性咨询 & 医生接触频率与 G1 匹配 & 非指导性、以确认为中心；无 CBT 会谈提示词、判断器、会谈评估或作业 & 完整记忆写入；保留咨询历史 & 非特异性因素与结构化 CBT 之对比 \\
G7 & 记忆写入消融 & 同 G1 & 与 G1 相同的完整 CBT & 主要通用记忆写入被阻断（事件、思想和聊天）；动态抑郁继续运行 & 新记忆形成的作用 \\
\bottomrule
\end{tabular}
\end{table}

%%====================================================================
%%  扩展数据（占位；按图表规划，终点表与主诉图热图降级至此）
%%====================================================================

\begin{figure}[h]
\centering
\fbox{\parbox[c][5cm][c]{0.9\textwidth}{\centering \emph{\ph{扩展数据图 1 占位}}\\治疗前后主诉图阶段性转移对比（G1 vs G7），按控制器和日志版本分层。}}
\par\smallskip
\parbox{0.9\textwidth}{\footnotesize\textbf{扩展数据图 1 |} 治疗前后的主诉图阶段性转移（G1 对 G7），按控制器和日志版本分层。有待同版本 G7 运行 \ph{待补充证据：同版本主诉图证据}。}
\end{figure}

\begin{figure}[h]
\centering
\fbox{\parbox[c][5cm][c]{0.9\textwidth}{\centering \emph{\ph{扩展数据图 2 占位}}\\测量信度：快照内重复生成分布、ICC(1,1)/ICC(1,$K$)、条目级一致性。}}
\par\smallskip
\parbox{0.9\textwidth}{\footnotesize\textbf{扩展数据图 2 |} 阶段性评估协议的测量信度：快照内重复生成分布、重复生成组内相关系数和条目级一致性。有待重复汇总计算 \ph{待补充证据}。}
\end{figure}

\begin{table}[h]
\parbox{0.9\textwidth}{\footnotesize\textbf{扩展数据表 1 |} 人设系统：临床情境案例库（九个人设，每个人设三个运行时可切换严重度档位，另有未分档默认配置）与八人设选择器集（18--66 岁；中度核心信念与可恢复活动；仅中度档）。严重度档位计数以当前检出的仓库为准；方法学摘要记录的是四档（36 种配置）布局，此分歧已标记待裁定。详细特征待目标批次核验。}
\par\smallskip
\centering
\footnotesize
\setlength{\tabcolsep}{4pt}
\begin{tabular}{@{}p{2.5cm}p{3.0cm}p{2.6cm}p{2.9cm}@{}}
\toprule
系统 & 人设 & 严重度档位 & 在实验中的角色 \\
\midrule
临床情境案例库 & 9 个情境案例（失业、社交羞耻、照护压力、创作受挫等） & 3 个运行时可切换档位（轻度、中度、重度）加未分档默认配置 & 严重度分层与试运行（如 0922 验证性试运行） \\
8 人设选择器集 & 8 个人设，18--66 岁（学校、职场、家庭、退休和独居情境） & 仅中度 & 自 0908 批次起的小镇实验 \\
\bottomrule
\end{tabular}
\end{table}

\begin{table}[h]
\parbox{0.9\textwidth}{\footnotesize\textbf{扩展数据表 2 |} 各组基线特征与终点结局（按图表规划从正文降级至此）。T0 和 S12 的 PHQ-9 与 BDI-II；S2--S10 的五条目短量表总分。所有数值待目标批次核验。}
\par\smallskip
\centering
\footnotesize
\begin{tabular}{@{}lccccc@{}}
\toprule
组别 & 描述 & 基线 PHQ-9 & 终点 PHQ-9 & $\Delta$PHQ-9 & $\Delta$BDI-II \\
\midrule
G1 & 结构化 CBT & \ph{X.X} & \ph{X.X} & \ph{X.X} & \ph{X.X} \\
G2 & 无干预 & \ph{X.X} & \ph{X.X} & \ph{X.X} & \ph{X.X} \\
G3 & 中性社会接触 & \ph{X.X} & \ph{X.X} & \ph{X.X} & \ph{X.X} \\
G4 & 咨询室场景 & \ph{X.X} & \ph{X.X} & \ph{X.X} & \ph{X.X} \\
G5 & 负面社会暴露 & \ph{X.X} & \ph{X.X} & \ph{X.X} & \ph{X.X} \\
G6 & 支持性咨询 & \ph{X.X} & \ph{X.X} & \ph{X.X} & \ph{X.X} \\
G7 & 记忆写入消融 & \ph{X.X} & \ph{X.X} & \ph{X.X} & \ph{X.X} \\
\bottomrule
\end{tabular}
\end{table}

\begin{table}[h]
\parbox{0.9\textwidth}{\footnotesize\textbf{扩展数据表 3 |} 仿真效应量与已发表基准的对比（心理治疗荟萃分析~\cite{cuijpers2013,hofmann2012,cuijpers2020}；对话式智能体递送 CBT 的证据~\cite{heinz2025,hang2026}）。\ph{待补充证据：单一锁定批次校准；如有偏离将透明报告，不做事后协议切换。}}
\par\smallskip
\centering
\footnotesize
\begin{tabular}{@{}lcc@{}}
\toprule
基准 & 已发表范围（$d$ 或 $g$） & 仿真的 G1 效应量 \\
\midrule
CBT 与对照（荟萃分析） & \ph{X.XX}--\ph{X.XX} & \ph{X.XX} \\
心理治疗与对照（荟萃分析） & \ph{X.XX}--\ph{X.XX} & \ph{X.XX} \\
AI 递送 CBT（RCT 与荟萃分析） & \ph{X.XX}--\ph{X.XX} & \ph{X.XX} \\
\bottomrule
\end{tabular}
\end{table}

%%====================================================================
%%  文末事项
%%====================================================================
\section*{数据可用性}

支持本研究发现的仿真输出数据——阶段性评估分数（T0 和 S12 的 PHQ-9、BDI-II；S2--S10 的五条目短量表）、主诉图阶段性转移轨迹、情绪向量轨迹和对话记录——可在 \ph{Repository}（\ph{DOI/URL}）获取。图 1--4、表~1 与各扩展数据项目的源数据随稿件一并提供。

\section*{代码可用性}

\ph{ModelName} 仿真平台——包括抑郁认知架构、干预与对照条件配置（G1--G7 及补充性 G9）、阶段性评估框架和批量实验编排脚本——公开发布于 \ph{GitHub URL}。仓库包含智能体配置文件、CBT 会谈提示词、居民聊天效价脚本、评估规则和批量实验脚本。平台以 Python 实现，采用 \ph{License} 许可证。

\section*{致谢}

我们感谢 \ph{姓名} 的 \ph{贡献}。本工作受 \ph{资助} 支持。我们声明使用大语言模型（DeepSeek、Qwen3-8B）作为研究基础设施；未经作者审阅，任何 LLM 输出均未纳入稿件正文。

\section*{作者贡献}

\ph{Author 1} 设计了抑郁认知架构，并实现了主诉图和情绪推断模块。\ph{Author 2} 开发了治疗干预管线和会谈控制系统。\ph{Author 3} 设计了实验框架并进行了统计分析。\ph{Author 4} 构思了本研究，指导了项目并撰写了稿件。所有作者审查并批准了最终稿件。

\section*{利益冲突}

作者声明无利益冲突。

%%====================================================================
%%  参考文献（编号式；v3 修订版参考文献库）
%%====================================================================
\renewcommand{\refname}{参考文献}
\bibliographystyle{unsrt}
\bibliography{nmh-references-v3}

\vspace{1em}
\hrule
\vspace{0.5em}
\noindent\small\textcolor{gray}{本文件为 nmh-manuscript-v3.tex 的中文对照检查版（生成日期 2026-09-27），用于作者审读修改论文结构、论证链和占位符位置；所有红色 \texttt{[...]} 为待填占位符。本文件不用于投稿，投稿请使用英文版 \texttt{sn-article-template/nmh-manuscript-v3.tex}（Springer Nature \texttt{sn-jnl.cls} 模板）。}

\end{document}
